% GENERATED FILE. Edit canonical lessons, then run npm run generate:books.
% canonical-source-hash: fnv1a64:ff78c66d0d221c4a
% canonical-lessons: FR-C344-raisonner, FR-C344-raisonnement, FR-C344-conclure, FR-C344-nuance, FR-C344-nuancer

\chapter{Raisonner}
\label{ch:fr-raisonner}
\hlchaptermodality{344}

\begin{chapteropening}
\textbf{By the end of this chapter:} \emph{I can reason, follow a line of reasoning, draw a conclusion and qualify what I said with a nuance.}

Complete the last lesson of chapter 344: I can reason, follow a line of reasoning, draw a conclusion and qualify what I said with a nuance.
\end{chapteropening}

\section[raisonner]{raisonner — to reason, to think logically}

\label{lesson:FR-C344-raisonner}



\begin{quote}
Before the new one: say the French for the criticism, then the French for to stress.
\end{quote}

\subsection*{You'll want to know: raisonner}
\textbf{raisonner} — \textquotedblleft{}to reason, to think logically\textquotedblright{}.

\begin{quote}
Il raisonne bien, mais il parle trop vite.
\end{quote}

\begin{quote}
\emph{He reasons well, but he speaks too fast.}
\end{quote}

From \emph{raison}. Not to be mixed up with \emph{résonner}, to resound.

\begin{cousinweb}
\emph{raisonner} is built on \emph{raison} \textquotedblleft{}reason\textquotedblright{}, which is inherited from Latin \emph{ratiō, ratiōnis} \textquotedblleft{}a reckoning, reason\textquotedblright{} (Lewis \& Short).
\end{cousinweb}

\begin{grammarlens}[title={What you've built}]
One more word for telling, explaining and describing at B1.
\end{grammarlens}

\subsection*{Your turn}
\begin{itemize}
\raggedright
  \item \emph{Say it:} \emph{raisonner}
  \item \emph{Say it:} \emph{raisonner}, once more
  \item \emph{Say it:} a sentence of your own with \emph{raisonner}
  \item \emph{Recall:} say the French for the criticism, then the French for to stress, then say \emph{raisonner} again
\end{itemize}

\begin{tcolorbox}[breakable,colback=teal!4,colframe=teal!35!black,title={Before you move on}]
What does raisonner mean? (\textquotedblleft{}To reason, to think logically\textquotedblright{}.) Say it once more.
\end{tcolorbox}
\section[le raisonnement]{le raisonnement — the reasoning}

\label{lesson:FR-C344-raisonnement}



\begin{quote}
Before the new one: say the French for to stress, then the French for to reason.
\end{quote}

\subsection*{You'll want to know: le raisonnement}
\textbf{le raisonnement} — \textquotedblleft{}the reasoning\textquotedblright{}.

\begin{quote}
Ton raisonnement est juste.
\end{quote}

\begin{quote}
\emph{Your reasoning is sound.}
\end{quote}

\emph{Un raisonnement} is one chain of thought. Masculine.

\begin{cousinweb}
\emph{raisonnement} is built from \emph{raisonner} and the suffix \emph{-ment} (EtymDB). The verb rests on \emph{raison}, from Latin \emph{ratiō} \textquotedblleft{}a reckoning, reason\textquotedblright{} (Lewis \& Short).
\end{cousinweb}

\begin{grammarlens}[title={What you've built}]
One more word for telling, explaining and describing at B1.
\end{grammarlens}

\subsection*{Your turn}
\begin{itemize}
\raggedright
  \item \emph{Say it:} \emph{le raisonnement}
  \item \emph{Say it:} \emph{le raisonnement}, once more
  \item \emph{Say it:} a sentence of your own with \emph{le raisonnement}
  \item \emph{Recall:} say the French for to stress, then the French for to reason, then say \emph{le raisonnement} again
\end{itemize}

\begin{tcolorbox}[breakable,colback=teal!4,colframe=teal!35!black,title={Before you move on}]
What does le raisonnement mean? (\textquotedblleft{}The reasoning\textquotedblright{}.) Say it once more.
\end{tcolorbox}
\section[conclure]{conclure — to conclude}

\label{lesson:FR-C344-conclure}



\begin{quote}
Before the new one: say the French for to reason, then the French for the reasoning.
\end{quote}

\subsection*{You'll want to know: conclure}
\textbf{conclure} — \textquotedblleft{}to conclude\textquotedblright{}.

\begin{quote}
Après la réunion, nous avons conclu qu'il fallait attendre.
\end{quote}

\begin{quote}
\emph{After the meeting, we concluded that we had to wait.}
\end{quote}

Past participle \emph{conclu}. \emph{Conclure un accord} — to close a deal.

\begin{cousinweb}
\emph{conclure} came into Old French as \emph{conclure}, a borrowing from Latin \emph{conclūdere} \textquotedblleft{}to shut up, to bring to an end\textquotedblright{} (EtymDB; Lewis \& Short), built on \emph{claudere} \textquotedblleft{}to shut\textquotedblright{}.
\end{cousinweb}

\begin{grammarlens}[title={What you've built}]
One more word for telling, explaining and describing at B1.
\end{grammarlens}

\subsection*{Your turn}
\begin{itemize}
\raggedright
  \item \emph{Say it:} \emph{conclure}
  \item \emph{Say it:} \emph{conclure}, once more
  \item \emph{Say it:} a sentence of your own with \emph{conclure}
  \item \emph{Recall:} say the French for to reason, then the French for the reasoning, then say \emph{conclure} again
\end{itemize}

\begin{tcolorbox}[breakable,colback=teal!4,colframe=teal!35!black,title={Before you move on}]
What does conclure mean? (\textquotedblleft{}To conclude\textquotedblright{}.) Say it once more.
\end{tcolorbox}
\section[la nuance]{la nuance — the subtle difference, the shade}

\label{lesson:FR-C344-nuance}



\begin{quote}
Before the new one: say the French for the reasoning, then the French for to conclude.
\end{quote}

\subsection*{You'll want to know: la nuance}
\textbf{la nuance} — \textquotedblleft{}the subtle difference, the shade\textquotedblright{}.

\begin{quote}
Il y a une petite nuance entre ces deux mots.
\end{quote}

\begin{quote}
\emph{There is a slight difference between these two words.}
\end{quote}

Of colours too: \emph{une nuance de bleu}.

\begin{cousinweb}
The origin of \emph{nuance} is usually traced to Old French \emph{nuer} \textquotedblleft{}to shade, like clouds\textquotedblright{}, from \emph{nue} \textquotedblleft{}cloud\textquotedblright{}, inherited from Latin \emph{nūbēs} \textquotedblleft{}cloud\textquotedblright{} (Lewis \& Short).
\end{cousinweb}

\begin{grammarlens}[title={What you've built}]
One more word for telling, explaining and describing at B1.
\end{grammarlens}

\subsection*{Your turn}
\begin{itemize}
\raggedright
  \item \emph{Say it:} \emph{la nuance}
  \item \emph{Say it:} \emph{la nuance}, once more
  \item \emph{Say it:} a sentence of your own with \emph{la nuance}
  \item \emph{Recall:} say the French for the reasoning, then the French for to conclude, then say \emph{la nuance} again
\end{itemize}

\begin{tcolorbox}[breakable,colback=teal!4,colframe=teal!35!black,title={Before you move on}]
What does la nuance mean? (\textquotedblleft{}The subtle difference, the shade\textquotedblright{}.) Say it once more.
\end{tcolorbox}
\section[nuancer]{nuancer — to qualify, to tone down (a view)}

\label{lesson:FR-C344-nuancer}



\begin{quote}
Before the new one: say the French for to conclude, then the French for the subtle difference.
\end{quote}

\subsection*{You'll want to know: nuancer}
\textbf{nuancer} — \textquotedblleft{}to qualify, to tone down (a view)\textquotedblright{}.

\begin{quote}
Je voudrais nuancer ce que j'ai dit hier.
\end{quote}

\begin{quote}
\emph{I would like to qualify what I said yesterday.}
\end{quote}

\emph{Nous nuançons}, with \emph{ç} before \emph{o}.

\begin{cousinweb}
\emph{nuancer} is built on \emph{nuance} \textquotedblleft{}shade\textquotedblright{}, whose origin is usually traced to Old French \emph{nue} \textquotedblleft{}cloud\textquotedblright{}, inherited from Latin \emph{nūbēs} \textquotedblleft{}cloud\textquotedblright{} (Lewis \& Short).
\end{cousinweb}

\begin{grammarlens}[title={What you've built}]
One more word for telling, explaining and describing at B1.
\end{grammarlens}

\subsection*{Your turn}
\begin{itemize}
\raggedright
  \item \emph{Say it:} \emph{nuancer}
  \item \emph{Say it:} \emph{nuancer}, once more
  \item \emph{Say it:} a sentence of your own with \emph{nuancer}
  \item \emph{Recall:} say the French for to conclude, then the French for the subtle difference, then say \emph{nuancer} again
\end{itemize}

\begin{tcolorbox}[breakable,colback=teal!4,colframe=teal!35!black,title={Before you move on}]
What does nuancer mean? (\textquotedblleft{}To qualify, to tone down (a view)\textquotedblright{}.) Say it once more.
\end{tcolorbox}
